%==========================================================================================
\section{Trade-offs}\label{sec:tradeoffs}

Each trade-off below is a result with its condition. Section~\ref{sec:dynamics} concerns the dynamics each
set-up generates. Section~\ref{sec:conclusions} concerns the conclusions each supports.

\subsection{Dynamics}\label{sec:dynamics}

\paragraph{The partial sale.}
\lap{At \(1.55\), \(1.60\) and \(1.65\) the certified equilibria have pure orders \((1,-v)\), with
\(v=\reg{lap_P3_v_r155}\), \reg{lap_P3_v_r160} and \reg{lap_P3_v_r165} (Proposition 3; computer-assisted). The
interior sale exists because the low class's marginal profit falls strictly with its order (Online Appendix
B.2). A root of its first-order condition is therefore a global optimum. The branch between the nodes is a
numerical diagnostic.}

\slo{On the unit menu a sale is all or nothing, and the only continuous margin is a mixture probability. The low
class abstains up to \(r_N=\reg{rN}\), mixes on \((r_N,r_S)\) and sells with probability one from
\(r_S=\reg{rS_B}\) (Section 4.3; computer-assisted). The mixture needs the investor rule IM. Under I0 the
equilibrium set is empty on \((r_N,r_S]\).}

The partial sale is therefore a property of the continuous order menu. Its condition is that the investor can
choose a size, with a density that moves smoothly with it.

\paragraph{Strictly rising entry at intermediate strengths.}
\lap{Entry rises strictly across the three nodes,
\(\reg{lap_P3_entry_r155}<\reg{lap_P3_entry_r160}<\reg{lap_P3_entry_r165}\), because the entry enclosures are
disjoint (Proposition 3; computer-assisted). The sale magnitude carries the rise: a larger \(v\) makes a sale
more informative and moves the threshold flow. The conditions are the continuous menu, a fixed presence of the
investor and an interval certificate at each node.}

\slo{With a fixed presence probability, entry rises strictly on \([r_A,r_F]=[\reg{rA},\reg{rF_B}]\), where class
\(H\) mixes (Section 4.3; computer-assisted). Entry is \reg{entry_rA_B} at \(r_A\), \reg{entry_mb_r145_B} at
\reg{str_r145} and \reg{entry_r1_B} from \(r_F\). The purchase rate \(s_+=b[\DT(r)/k-2]\) adjusts so that a
purchase that recruits earns exactly \(k\). The rise needs rule IM, and entry is flat on \([r_F,r_C]\); at
\(1.55\) to \(1.65\) it equals \reg{entry_r1_B}.}

\slo{With research, entry rises strictly on \([\reg{res_rR},\reg{res_rcap})\), which contains the three
strengths: \(\reg{res_entry_r155}<\reg{res_entry_r160}<\reg{res_entry_r165}\) (Section 5; slope analytical,
interval computer-assisted). The research intensity carries the rise. It solves \(2a\lambda=\DT(r)m_\lambda-k\),
and class \(H\) buys with probability one after a success. The conditions are a quadratic research cost with
\(a\) in \((\reg{res_a_lo},\reg{res_a_up})\) and a cap in \((\reg{res_cap_lo},\reg{res_cap_hi})\). Zero research
is an equilibrium at every strength on the branch.}

Each set-up therefore has a strict rise at intermediate strengths. A different object carries it in each case: a
sale size, a mixture probability or a research intensity. The slot's fixed-presence rise sits at lower strengths
than the Laplace nodes. Its research rise sits at the Laplace nodes with lower entry levels.

\paragraph{Coexistence.}
\lap{No trade is an equilibrium exactly on \(r\le r_N=\reg{lap_rN}\) and is unique below
\(\fr(k)=\reg{lap_r_unique_zero}\). Full orders are unique above \(r_U=\reg{lap_rU}\) (Proposition A.4;
analytical). Between the two bounds the manuscript certifies three informative equilibria, each beside no trade,
and proves no uniqueness. The equilibrium set there is open.}

\slo{No trade and an informative equilibrium coexist exactly on \([r_A,r_N]\), and the set at each strength is
listed (Section 4.3; computer-assisted). The high class uses one of five forms: no trade, a mixed purchase that
recruits, a full purchase, mixed orders at the entry floor, or the threshold pair. Under rule CM the threshold
pair is also an equilibrium on \((r_A,r_J]\), \(r_J=\reg{rJ}\), with entry that falls strictly from
\reg{entry_rA_B} to \(\rho\). At \(r_C\) under CM, entry after a purchase can take any value in \([\rho,1]\).}

The slot's completeness follows from a finite support enumeration. Each class has three orders, and each gross
return falls strictly with the class's own rate. The Laplace set-up has no such enumeration on a continuum.

\paragraph{Research.}
Only the slot models research. The investor chooses an intensity \(z\in[0,\bar\lambda]\) at cost \((a/2)z^2\)
before it learns the class, and the intensity is the presence probability. Zero research is an equilibrium
exactly when \(\rho\DT\le2k\) at an uninformative price; this condition does not involve \(a\). The marginal
cost rises and the research value falls, so they meet once, and the research response moves smoothly with
strength. The benchmark cap sits \reg{margin_lambda_hi_B} below the top of the admissible strip, the same
boundary as the return margin at \(r_2\) in another coordinate. The equilibrium set with research on
\((r_N,r_U^R)=(\reg{res_rN},\reg{res_rUR})\) is open.

\paragraph{Timing.}
Both set-ups price one trade before the challenger's decision, and neither models the choice of a trading
date. The Laplace flow is the aggregate net demand of the decision interval. The slot is one public order
opportunity at a known date. This is a stronger institutional assumption with a sharper empirical counterpart:
the first prespecified opportunity, with a recorded zero when no order arrives. An exploratory extension with
several slots, outside both manuscripts, keeps the posterior bound and the return to the floor at \(r_2\) for
every finite number of slots; its complete equilibrium list is open. In both set-ups the ceiling \(r_C\) is a
tie event with positive probability, so entry is discontinuous in strength there and the challenger's tie rule
decides entry at \(r_C\) itself.

\subsection{Conclusions each set-up supports}\label{sec:conclusions}

\paragraph{The fixed-information control: falls, or does not rise.}
\lap{Hold full orders fixed and let the challenger decide again. Entry falls strictly from
\reg{lap_frozen_r0_entry} at \(r_0\) to \reg{lap_frozen_r1_entry} at \(r_1\) (Section 4.2; closed forms). The
threshold flow \(x^*\) moves continuously with strength, so the crossing probability has a density and the fall
is strict.}

\slo{With the same control, entry is unchanged at \reg{p4_frozen_r0_entry} on \((1,r_C]\) and equals \(\rho\)
above \(r_C=\reg{rC_B}\) (Section 4.2; analytical, values computer-assisted). The frozen posteriors are three
atoms \(m<1/2<M\), and the threshold \(\tau(r)\) exceeds \(5/8\) for \(r>1\). Only the top atom can recruit, so
entry is a step function of strength. The crossing probability is \reg{p4_control_crossing_r0_r1} from \(r_0\)
to \(r_1\) and \reg{p4_control_crossing_r1_r2} from \(r_1\) to \(r_2\).}

The condition is the number of prices. With three prices entry moves in steps, so a strict fall at fixed orders
needs the threshold to cross an atom. The weak statement ``does not rise'' holds in both set-ups.

\paragraph{Reserve deterrence: strict, or flat.}
\lap{The reserve \reg{lap_reserve_excl_p} excludes the low class. There, preparation falls from
\reg{lap_reserve_excl_r0_entry} at \(r_0\) to \reg{lap_reserve_excl_r1_entry} at \(r_1\), with unique
continuations at both nodes (Table 5; analytical).}

\slo{At the reserve \reg{rsv_points_excl}, entry is \reg{p2_reserve_points_exclude_r0_entry} at both strengths:
full orders are unique at both, and only a purchase recruits at both. Strict deterrence returns at the reserve
\reg{rsv_deter}, where entry falls from \reg{p2_reserve_points_deter_r0_entry} to
\reg{p2_reserve_points_deter_r1_entry} (Section 7.3; analytical, values computer-assisted). This reserve
exceeds \(r_0\) and excludes the weak incumbent. Strict deterrence after the low class is excluded holds for
reserves between \(p_H\approx\reg{p2_reserve_ph}\) and \(h-c_H/M\approx\reg{p2_reserve_weak_top}\).}

The condition is the same step structure. Strict deterrence in the slot needs the stronger incumbent to push the
top-posterior profit below the high cost.

\paragraph{Uniqueness and completeness.}
The Laplace set-up proves uniqueness against a continuum of deviations at each benchmark strength and leaves
the set between \(r_N\) and \(r_U\) open. The slot proves uniqueness against three orders and lists the whole
set along the path, with the dependence on the two tie rules in one table (Section 4.3). Three statements depend
on the rules. The strict rise on \([r_A,r_F]\) needs IM; the threshold pair and the entry interval at \(r_C\)
exist only under CM; under I0 the set is empty on \((r_N,r_S]\). Proposition 2 and the research results at the
six benchmark strengths hold under all four pairs.

\paragraph{The admissible region.}
Both manuscripts have the same thin boundary at \(r_2\). The Laplace manuscript asserts that its region is open
and prints prior windows; it does not print the part (iii) margin \(c_H-B_{r_2}(M)=\reg{lap_return_cost}\). The
slot manuscript states the region as an exact band in \(u\) with a table of margins. In the slot the thin margin
appears in five coordinates that measure one distance. They are \reg{margin_return_post_B} in the top
posterior, \reg{margin_return_B} in cost units, \reg{margin_u_hi_B} in \(u\), \reg{margin_lambda_hi_B} in
\(\lambda\) at the benchmark \(\beta\), and \reg{gap_r2_rC_B} in strength. Section~\ref{sec:worked} gives the
windows in \(k\) and in the prior, which differ between the set-ups.

\paragraph{The empirical mapping.}
\lap{Section 7 states three predictions. A stronger prepared bidder raises the probability of further
preparation when the price precedes the decision, and weakly lowers it when it does not. A buyer that prepares
after a favorable return wins more often. Against a strong prepared bidder, preparation returns to its floor
while the price stays informative. The flow is aggregate net demand, so the model's largest price gap is not a
unit spread, and no observable is derived from the flow.}

\slo{The same three predictions hold (Section 8, in preparation), and the slot adds derived observables. The
conditional execution spread is \(S=\rho\DT\vartheta+(1-\rho)J\,A(M)\), with \(\vartheta=u/(2+u)\) the informed
share of nonzero orders. It is \reg{ba_r1_B} at \(r_1\), and \reg{ba_share_r1_B} of it is the entry component;
the entry component is zero above \(r_C\). Sell impact and volume identify \((\vartheta,\lambda,\beta)\) under
full orders, while counts alone leave the informed share in
\([\reg{theta_min_counts_B},\reg{theta_max_full_eq_r1_B})\). Orders predict the winner among prepared
challengers above \(r_C\).}

The slot's formulas need an observation protocol: the first prespecified public opportunity, with a recorded
zero. A sample of first executed trades cannot identify the neutral price effect. The formulas are analytical,
the values computer-assisted, and the data counterpart of each protocol is open.

\paragraph{Assumptions each set-up adds.}
The Laplace set-up assumes a continuous order, a Laplace law with \(b>1\) and the score bound; Proposition A.5
extends it to every bounded-score density. The slot assumes a unit menu, a presence probability and an exogenous
liquidity motive. An uninformed trader with the model's own trade payoff would not fill the slot. With research
the slot adds a quadratic cost and a cap inside the admissible strip. Neither set-up needs an off-path belief.

